%% ma: L10-B15-0053
%% lop: 10
%% chuong: 6
%% bai: 15
%% dang: 10.15.6
%% loai: TL
%% muc: VD
%% dap_an: "-"
%% nguon: "Ảnh giáo viên gửi 21/11/2026: Chuyên đề VI – Toán 10 – Chương VI (Hàm số – đồ thị và ứng dụng), tr. 199–203, Dạng 6 Câu 3"
%% kien_thuc: "Lớp 10 Bài 15 (hàm số trong thực tế); hằng đẳng thức (lớp 8)"
%% trang_thai: chinh-thuc
%% ghi_chu: "Dạng do giáo viên cung cấp (21/11/2026)."
\begin{ex}
Một cửa hàng buôn giày nhập một đôi với giá là $40$ USD. Cửa hàng ước tính rằng nếu đôi giày được bán với giá $x$ USD thì mỗi tháng khách hàng sẽ mua $(120-x)$ đôi. Hỏi cửa hàng bán một đôi giày giá bao nhiêu thì thu được nhiều lãi nhất?
\loigiai{Lãi mỗi tháng: $L(x)=(x-40)(120-x)=-x^2+160x-4800=1600-(x-80)^2\le1600$ (với $40\le x\le120$). Dấu bằng khi $x=80$. Vậy bán giá $80$ USD một đôi, lãi lớn nhất $1600$ USD mỗi tháng.}
\end{ex}
